% Biodata anggota pengganti (format CV mini proposal) + KTP dengan NIK disensor.
% Kompilasi: latexmk -pdf biodata-budi.tex
\documentclass[11pt,a4paper]{article}
\usepackage[a4paper,margin=2.54cm]{geometry}
\usepackage[T1]{fontenc}
\usepackage[utf8]{inputenc}
\usepackage[scaled=0.92]{helvet}
\renewcommand{\familydefault}{\sfdefault}
\usepackage{microtype}
\usepackage{setspace}
\setstretch{1.15}
\setlength{\parindent}{0pt}
\setlength{\parskip}{6pt}
\usepackage[table]{xcolor}
\usepackage{graphicx}
\usepackage{tabularx}
\usepackage{array}
\usepackage{titlesec}
\usepackage{tikz}
\usepackage[hidelinks]{hyperref}
\usepackage{xurl}
\urlstyle{same}
\frenchspacing

\definecolor{hdrgray}{HTML}{F2F2F2}
\newcommand{\hdr}[1]{\cellcolor{hdrgray}\textbf{#1}}
\newcommand{\todo}[1]{\textcolor{red}{\textbf{[#1]}}}
\newcolumntype{L}[1]{>{\raggedright\arraybackslash}p{#1}}
\newcolumntype{Y}{>{\raggedright\arraybackslash}X}
\renewcommand{\arraystretch}{1.25}
\titleformat{\section}{\large\bfseries\raggedright}{}{0em}{}
\titleformat{\subsection}{\normalsize\bfseries\raggedright}{}{0em}{}
\pagestyle{empty}

% KTP dengan NIK sudah disensor dan diberi watermark langsung di gambarnya (ktp-budi.png).
\newcommand{\ktpfile}{ktp-budi}

\begin{document}

\begin{center}
  {\Large\bfseries Biodata Anggota Tim}\\[2pt]
  Hackathon Deeptech 2026: Indonesia Chip Hackathon Challenge 2026\\[2pt]
  Tim NaS --- SeBatt: Secure Battery Identity Chip
\end{center}

\section{C. CV Anggota}

\subsection{Anggota 2 --- Muhammad Budi Setiawan}

\begin{tabularx}{\linewidth}{|L{4.6cm}|Y|}
  \hline
  \hdr{Item} & \hdr{Isian} \\ \hline
  Nama lengkap & Muhammad Budi Setiawan \\ \hline
  NIM & 22/505064/TK/55254 \\ \hline
  Program studi, angkatan/ Instansi/Perusahaan & Teknologi Informasi, 2022, Universitas Gadjah Mada \\ \hline
  Keahlian teknis & \textit{Machine learning} untuk data baterai (LSTM, \textit{reinforcement learning}), estimasi \textit{State of Charge}, Python \\ \hline
  Pengalaman proyek, riset, atau magang & Skripsi: Estimasi \textit{State of Charge} Baterai Li-Ion Menggunakan LSTM dengan Koreksi Residual \textit{Reinforcement Learning}; Data Engineering Intern, PT Packet System Indonesia (Januari--Februari 2025), membangun \textit{pipeline} data \textit{end-to-end} dari Azure Log Analytics ke Data Lake Microsoft Fabric hingga basis data SQL; Pengembangan \textit{frontend} web Pionir Kesatria FT UGM, Technocorner, dan KPU FT UGM dengan Next.js \\ \hline
  Prestasi atau lomba & - \\ \hline
  Portofolio & \url{https://www.linkedin.com/in/mbudis/} \\ \hline
\end{tabularx}

\newpage
\section{Kartu Tanda Penduduk}

\IfFileExists{\ktpfile.jpg}{\def\ktpok{1}}{\IfFileExists{\ktpfile.png}{\def\ktpok{1}}{\def\ktpok{0}}}
\if1\ktpok
\begin{center}
\includegraphics[width=0.9\linewidth]{\ktpfile}
\end{center}
\else
\todo{File KTP belum ada. Simpan sebagai anggota/ktp-budi.jpg lalu kompilasi ulang.}
\fi

\end{document}
